\section{层次化动机与任务拆解}

\subsection{长视野的典型陷阱}

单层策略面对稀疏奖励、组合爆炸与非平稳环境时常常陷入“微调-过拟合”的先验，Jason Punchline 是不放弃结构化视角。通过添加层次，我们重新定义动作空间的粒度：每个宏 action 负责若干步的子任务，降低信用分配压力。

\begin{knowledgebox}{长视野结构化分解的三步}
\begin{enumerate}
    \item 识别关键 check-point（example：厨房任务的“取出锅”）；
    \item 提取每个 check-point 所需的 subgoal；
    \item 设计子策略完成该 subgoal，并在完成后反馈全局 reward。
\end{enumerate}
\end{knowledgebox}

\begin{warningbox}{忽略分解的代价}
\begin{itemize}
    \item 在平坦策略中，长轨迹上任意一次失败都需要重新探索；
    \item Credit assignment 模型容易被噪音启动子任务，导致训练震荡；
    \item 没有分层的代理难以压缩上下文，面对部分可观测环境时崩溃。
\end{itemize}
\end{warningbox}

\subsection{层次化与归纳偏置}

层次化在某种程度上也是归纳偏置——我们告诉模型“先完成一系列有意义的中间目标再接尾声”。正确地注入这样的偏置可以显著降低 sample complexity；但如果偏置不对（例如固定 skill 数量、缺少通用 stop condition），会形成 training bottleneck。

\begin{table}[H]
\centering
\begin{tabular}{p{4cm}p{5cm}p{3cm}}
\toprule
维度 & 传统 RL & 层次化 RL \\
\midrule
决策粒度 & 原子动作 & 子策略/宏动作 \\
样本效率 & 低，需完整 rollout & 高，子策略可复用 \\
探索策略 & 纯探索 & 构造 subgoal 降低 search \\
信用分配 & 全局回报 & 子策略附加 reward \\
\bottomrule
\end{tabular}
\caption{传统 RL 与层次化 RL 任务拆解的对比。}
\end{table}

\subsection{工业场景中的分层模板}

在实际产品（如厨房机器人、自动化测试平台）中，层次化模板通常包括：1) 感知/解析模块（抽取 subgoal），2) Strategy Planner（选 option/high-level plan），3) Execution Layer（skill + verifiers），4) Recovery + Logging（incident grade-book + prompt log）。每层都需把 input/output 明文记录到 `notes/ops/`，以便 on-call 团队快速对照。

\begin{knowledgebox}{常见工业模板}
\begin{itemize}
    \item Perception：state filtering + object detection；
    \item Planner：high-level policy 结合 options；
    \item Executor：skill library + environment interaction；
    \item Recovery：failure flag + fallback skill。
\end{itemize}
\end{knowledgebox}

\subsection{本章小结}

动机部分强调“结构化”不是 add-on，而是解决信用分配和 exploration 的核心；后续的 Options、goal-conditioned 和模仿部分将围绕这两条主线展开。

